\documentclass[11pt]{article}
\usepackage[T1]{fontenc}
\usepackage{lmodern}
\usepackage[ngerman]{babel}
\usepackage[a4paper,margin=2cm]{geometry}
\usepackage{amsmath,amssymb,graphicx,xcolor,enumitem,multicol}
\pagestyle{empty}
\setlength{\parindent}{0pt}
\setlist[enumerate,1]{label=(\arabic*), leftmargin=2.2em, itemsep=10pt}
\newcommand\kaestchen[1]{\makebox[1.4em][l]{$\square$}\textbf{#1}\enspace}
\begin{document}
{\sffamily\bfseries\Large BLF-Training Mathematik – Multiple Choice}\par\smallskip
Name: \rule{7cm}{0.4pt}\hfill Datum: \rule{3cm}{0.4pt}\par\medskip
{\small Kreuze die richtige Antwort an. Bei Aufgaben mit \textbf{(M)} können mehrere Antworten richtig sein.}
\par\medskip\hrule\medskip
\begin{enumerate}
\item Welche Gleichung gehört zu einer linearen Funktion mit dem Anstieg $2$, deren Graph die $y$-Achse im Punkt $(0 \mid -3)$ schneidet?\par\nopagebreak
\begin{multicols}{2}
\kaestchen{A}$f(x) = x^2 - 3$\par
\kaestchen{B}$f(x) = 2x - 3$\par
\kaestchen{C}$f(x) = 2x + 3$\par
\kaestchen{D}$f(x) = -3x + 2$\par
\end{multicols}
\item Gib die Lösungen der Gleichung $x^2 - 5x + 6 = 0$ an.\par\nopagebreak
\begin{multicols}{2}
\kaestchen{A}keine reelle Lösung\par
\kaestchen{B}$x_1 = -2;\ x_2 = -3$\par
\kaestchen{C}$x_1 = 2;\ x_2 = 3$\par
\kaestchen{D}$x_1 = 1;\ x_2 = 6$\par
\end{multicols}
\item Ein Artikel kostet 80 €. Der Preis wird um 15 \% gesenkt. Wie viel kostet er danach?\par\nopagebreak
\begin{multicols}{2}
\kaestchen{A}68 €\par
\kaestchen{B}72 €\par
\kaestchen{C}65 €\par
\kaestchen{D}92 €\par
\end{multicols}
\item Ein Zylinder hat den Radius $r = 3\,\text{cm}$ und die Höhe $h = 10\,\text{cm}$. Wie groß ist sein Volumen (gerundet)?\par\nopagebreak
\begin{multicols}{2}
\kaestchen{A}$188\,\text{cm}^3$\par
\kaestchen{B}$90\,\text{cm}^3$\par
\kaestchen{C}$283\,\text{cm}^3$\par
\kaestchen{D}$94\,\text{cm}^3$\par
\end{multicols}
\item Gegeben ist $f(x) = -(x-1)^2 + 4$. Welche Aussagen sind wahr? \textbf{(M)}\par\nopagebreak
\kaestchen{A}Der Graph schneidet die $y$-Achse im Punkt $(0 \mid 4)$.\par
\kaestchen{B}Die Nullstellen sind $x_1 = -1$ und $x_2 = 3$.\par
\kaestchen{C}Die Parabel ist nach unten geöffnet.\par
\kaestchen{D}Der Scheitelpunkt ist $S(1 \mid 4)$.\par
\item Welchen Wert hat $\sin 30^\circ$?\par\nopagebreak
\begin{multicols}{2}
\kaestchen{A}$\frac{\sqrt2}{2}$\par
\kaestchen{B}$\frac{\sqrt3}{2}$\par
\kaestchen{C}$1$\par
\kaestchen{D}$0{,}5$\par
\end{multicols}
\item Es wird zweimal gewürfelt. Wie groß ist die Wahrscheinlichkeit für die Augensumme 12?\par\nopagebreak
\begin{multicols}{2}
\kaestchen{A}$\frac{1}{12}$\par
\kaestchen{B}$\frac{1}{36}$\par
\kaestchen{C}$\frac{2}{36}$\par
\kaestchen{D}$\frac{1}{6}$\par
\end{multicols}
\item Ein rechtwinkliges Dreieck hat die Katheten $6\,\text{cm}$ und $8\,\text{cm}$. Wie lang ist die Hypotenuse?\par\nopagebreak
\begin{multicols}{2}
\kaestchen{A}$48\,\text{cm}$\par
\kaestchen{B}$7\,\text{cm}$\par
\kaestchen{C}$14\,\text{cm}$\par
\kaestchen{D}$10\,\text{cm}$\par
\end{multicols}
\item Eine Bakterienkultur verdoppelt sich jede Stunde. Am Anfang sind es 100 Bakterien. Wie viele sind es nach 5 Stunden?\par\nopagebreak
\begin{multicols}{2}
\kaestchen{A}1000\par
\kaestchen{B}500\par
\kaestchen{C}1600\par
\kaestchen{D}3200\par
\end{multicols}
\item Welche Darstellung ist gleich $0{,}00042$?\par\nopagebreak
\begin{multicols}{2}
\kaestchen{A}$4{,}2 \cdot 10^{4}$\par
\kaestchen{B}$4{,}2 \cdot 10^{-4}$\par
\kaestchen{C}$4{,}2 \cdot 10^{-3}$\par
\kaestchen{D}$4{,}2 \cdot 10^{-5}$\par
\end{multicols}
\end{enumerate}
\clearpage\scriptsize\raggedright
{\sffamily\bfseries\normalsize Lösungen: BLF-Training Mathematik – Multiple Choice}\par\smallskip\hrule\medskip
\begin{multicols}{2}\begin{enumerate}
\item $\boxed{\text{B}}$ -- $f(x) = mx + n$ mit $m = 2$ und $n = -3$.
\item $\boxed{\text{C}}$ -- p-q-Formel mit $p = -5$, $q = 6$: $x = 2{,}5 \pm 0{,}5$.
\item $\boxed{\text{A}}$ -- $80 \cdot 0{,}85 = 68$, also 68 €.
\item $\boxed{\text{C}}$ -- $V = \pi r^2 h = \pi \cdot 9 \cdot 10\,\text{cm}^3 \approx 282{,}7\,\text{cm}^3$
\item $\boxed{\text{B, C, D}}$ -- $f(0) = -1 + 4 = 3$, also Schnittpunkt $(0 \mid 3)$.
\item $\boxed{\text{D}}$
\item $\boxed{\text{B}}$ -- Nur das Ergebnis $(6;6)$ von 36 gleich wahrscheinlichen Ergebnissen.
\item $\boxed{\text{D}}$ -- $\sqrt{36 + 64} = \sqrt{100} = 10$
\item $\boxed{\text{D}}$ -- $100 \cdot 2^5 = 3200$
\item $\boxed{\text{B}}$
\end{enumerate}\end{multicols}
\end{document}